\begin{proof}[Proof of \cref{obj:lem:prior-reduction}]
\begin{enumerate}
\item Choose the universal constant \(c_{\mathrm{floor}}>0\) supplied by
\cref{obj:lem:parametric-floor-infrastructure}, so that
\[
\mathfrak R^*_{n,d,q}\ge \frac{c_{\mathrm{floor}}}{n}
\]
throughout the stated parameter range. Define
\[
c_{\mathrm{sep,pt}}
=\frac{13}{32}c_{\beta_0}^{\,2},
\qquad
c=\frac12\min\{c_{\mathrm{floor}},c_{\mathrm{sep,pt}}\}.
\]
Since \(\beta_0=1/16\in(0,1/4)\), the selected constant in
\cref{obj:thm:normalized-converse} satisfies \(c_{\beta_0}>0\).
Consequently,
\[
c_{\mathrm{sep,pt}}>0,
\qquad c>0.
\]

Fix \(\alpha\in(0,1/2)\). Choose \(c_{0,\alpha}>0\) from
\cref{obj:lem:parametric-floor-infrastructure}, giving
\[
\mathfrak L^*_{n,d,q,\alpha}
\ge \frac{c_{0,\alpha}}{\sqrt n}
\]
throughout the same parameter range. Define
\[
\beta_\alpha=\frac{1-2\alpha}{8},
\qquad
c_{\mathrm{sep,len},\alpha}
=\frac54c_{\beta_\alpha}(1-2\alpha),
\qquad
c_\alpha
=\frac12\min\{c_{0,\alpha},c_{\mathrm{sep,len},\alpha}\}.
\]
The inequalities
\[
0<1-2\alpha<1,
\qquad
0<\beta_\alpha<\frac18<\frac14
\]
ensure \(c_{\beta_\alpha}>0\), hence
\[
c_{\mathrm{sep,len},\alpha}>0,
\qquad c_\alpha>0.
\]
The constants \(c,c_{\mathrm{floor}}\) are universal, and \(c_\alpha,c_{0,\alpha}\)
depend only on \(\alpha\).
% lean: CausalSmith.Stat.MarNearcompleteFrontier.prior_reduction

\item Fix admissible \(n,d,q\). The scale is nonnegative because
\[
1-q\ge0,\qquad
\ell_n=\log(e+n)>0,\qquad
\frac{d}{n\ell_n}\ge0,
\qquad
g_{n,d,q}=(1-q)\min\left\{1,\frac{d}{n\ell_n}\right\}\ge0.
\]
Suppose \(g_{n,d,q}^{\,2}\ge n^{-1}\), and define
\[
a_{\mathrm{pt}}=c_{\beta_0}g_{n,d,q}\ge0.
\]
By \cref{obj:thm:normalized-converse}, there are probability priors
\(\Pi_{-,\mathrm{pt}},\Pi_{+,\mathrm{pt}}\) and a center
\(m_{0,\mathrm{pt}}\) satisfying
\[
\Pi_{\sigma,\mathrm{pt}}\bigl(\mathcal M(d,q)\bigr)=1,
\qquad
m_{0,\mathrm{pt}}\in\mathbb R,
\qquad \sigma\in\{-,+\}.
\]
Define their sample-mixture laws by
\[
Q_{\sigma,\mathrm{pt}}^{(n)}
=\int_{\mathcal M(d,q)}P_O^{\otimes n}\,
  \Pi_{\sigma,\mathrm{pt}}(dP),
\qquad \sigma\in\{-,+\}.
\]
The supplied bounds are
\[
\operatorname{TV}\!\left(
Q_{-,\mathrm{pt}}^{(n)},Q_{+,\mathrm{pt}}^{(n)}
\right)\le\beta_0,
\]
\[
\Pi_{-,\mathrm{pt}}\{\tau(P)\le m_{0,\mathrm{pt}}-a_{\mathrm{pt}}\}
\ge1-\beta_0,
\qquad
\Pi_{+,\mathrm{pt}}\{m_{0,\mathrm{pt}}+a_{\mathrm{pt}}\le\tau(P)\}
\ge1-\beta_0.
\]

For an arbitrary estimator \(T\in\mathcal T_n\), define
\[
E_{\mathrm{pt}}
=\{o\in O^n:m_{0,\mathrm{pt}}\le T(o)\},
\]
\[
\mathcal B_{\sigma,\mathrm{pt}}(T)
=\int_{\mathcal M(d,q)}
  \mathbb E_P\!\left[(T(O^n)-\tau(P))^2\right]\,
  \Pi_{\sigma,\mathrm{pt}}(dP),
\qquad \sigma\in\{-,+\},
\]
and the target-failure indicators
\[
D_{-,\mathrm{pt}}(P)
=\mathbf1\{\tau(P)>m_{0,\mathrm{pt}}-a_{\mathrm{pt}}\},
\qquad
D_{+,\mathrm{pt}}(P)
=\mathbf1\{\tau(P)<m_{0,\mathrm{pt}}+a_{\mathrm{pt}}\}.
\]
On the left target region, \(o\in E_{\mathrm{pt}}\) gives
\[
T(o)-\tau(P)\ge a_{\mathrm{pt}},
\qquad
(T(o)-\tau(P))^2\ge a_{\mathrm{pt}}^2.
\]
On the right target region, \(o\in E_{\mathrm{pt}}^c\) gives
\[
\tau(P)-T(o)\ge a_{\mathrm{pt}},
\qquad
(T(o)-\tau(P))^2\ge a_{\mathrm{pt}}^2.
\]
Integrating these pointwise bounds over the corresponding sample events
gives the good-target risk bounds. On each failure region, the sample-event
probability is at most one and squared loss is nonnegative. Thus, for every
\(P\in\mathcal M(d,q)\),
\[
\begin{aligned}
a_{\mathrm{pt}}^2
\left[P_O^{\otimes n}(E_{\mathrm{pt}})-D_{-,\mathrm{pt}}(P)\right]
&\le \mathbb E_P[(T(O^n)-\tau(P))^2],\\
a_{\mathrm{pt}}^2
\left[P_O^{\otimes n}(E_{\mathrm{pt}}^c)-D_{+,\mathrm{pt}}(P)\right]
&\le \mathbb E_P[(T(O^n)-\tau(P))^2].
\end{aligned}
\]
The prior failure probabilities satisfy
\[
\int D_{\sigma,\mathrm{pt}}(P)\,
  \Pi_{\sigma,\mathrm{pt}}(dP)\le\beta_0,
\qquad \sigma\in\{-,+\}.
\]
Integration over the two priors therefore yields
\[
\begin{aligned}
\mathcal B_{-,\mathrm{pt}}(T)
&\ge a_{\mathrm{pt}}^2
 \left[Q_{-,\mathrm{pt}}^{(n)}(E_{\mathrm{pt}})-\beta_0\right],\\
\mathcal B_{+,\mathrm{pt}}(T)
&\ge a_{\mathrm{pt}}^2
 \left[Q_{+,\mathrm{pt}}^{(n)}(E_{\mathrm{pt}}^c)-\beta_0\right].
\end{aligned}
\]
Total variation gives
\[
\begin{aligned}
Q_{-,\mathrm{pt}}^{(n)}(E_{\mathrm{pt}})
+Q_{+,\mathrm{pt}}^{(n)}(E_{\mathrm{pt}}^c)
&=1+Q_{-,\mathrm{pt}}^{(n)}(E_{\mathrm{pt}})
     -Q_{+,\mathrm{pt}}^{(n)}(E_{\mathrm{pt}})\\
&\ge1-\beta_0.
\end{aligned}
\]
Adding the two risk bounds and substituting \(\beta_0=1/16\), we obtain
\[
\mathcal B_{-,\mathrm{pt}}(T)+\mathcal B_{+,\mathrm{pt}}(T)
\ge (1-3\beta_0)a_{\mathrm{pt}}^2
=\frac{13}{16}c_{\beta_0}^{\,2}g_{n,d,q}^{\,2}.
\]
Each prior average is at most the worst-law risk, since its prior is
supported on \(\mathcal M(d,q)\). Hence
\[
2\sup_{P\in\mathcal M(d,q)}
  \mathbb E_P[(T(O^n)-\tau(P))^2]
\ge
\mathcal B_{-,\mathrm{pt}}(T)+\mathcal B_{+,\mathrm{pt}}(T)
\ge 2c_{\mathrm{sep,pt}}g_{n,d,q}^{\,2}.
\]
The estimator and target both lie in \([-1,1]\), so the squared loss lies
in \([0,4]\); these bounds justify the prior integrations and the finite
worst-law risks. The estimator class contains the constant zero estimator.
Taking its infimum as in \cref{obj:def:point-risk} proves
\[
\mathfrak R^*_{n,d,q}
\ge c_{\mathrm{sep,pt}}g_{n,d,q}^{\,2}
=\frac{13}{32}c_{\beta_0}^{\,2}g_{n,d,q}^{\,2}.
\]
% lean: CausalSmith.Stat.MarNearcompleteFrontier.priorPointRisk_separation

\item Continue under \(g_{n,d,q}^{\,2}\ge n^{-1}\), and define
\[
a_{\mathrm{len}}=c_{\beta_\alpha}g_{n,d,q}\ge0.
\]
Apply \cref{obj:thm:normalized-converse} at \(\beta_\alpha\).
It supplies probability priors and a center with
\[
\Pi_{\sigma,\mathrm{len}}\bigl(\mathcal M(d,q)\bigr)=1,
\qquad
m_{0,\mathrm{len}}\in\mathbb R,
\qquad \sigma\in\{-,+\}.
\]
Define
\[
Q_{\sigma,\mathrm{len}}^{(n)}
=\int_{\mathcal M(d,q)}P_O^{\otimes n}\,
  \Pi_{\sigma,\mathrm{len}}(dP),
\qquad \sigma\in\{-,+\}.
\]
Then
\[
\operatorname{TV}\!\left(
Q_{-,\mathrm{len}}^{(n)},Q_{+,\mathrm{len}}^{(n)}
\right)\le\beta_\alpha,
\]
\[
\Pi_{-,\mathrm{len}}\{\tau(P)\le m_{0,\mathrm{len}}-a_{\mathrm{len}}\}
\ge1-\beta_\alpha,
\qquad
\Pi_{+,\mathrm{len}}\{m_{0,\mathrm{len}}+a_{\mathrm{len}}\le\tau(P)\}
\ge1-\beta_\alpha.
\]

For an arbitrary honest interval, write
\[
I\in\mathcal C_\alpha(n,d,q),
\qquad
I(o)=[L_I(o),U_I(o)],
\qquad o\in O^n,
\]
and define the endpoint events
\[
E_{\mathrm L,\alpha}
=\{o\in O^n:L_I(o)\le m_{0,\mathrm{len}}-a_{\mathrm{len}}\},
\qquad
E_{\mathrm R,\alpha}
=\{o\in O^n:m_{0,\mathrm{len}}+a_{\mathrm{len}}\le U_I(o)\}.
\]
Also define
\[
D_{-,\mathrm{len}}(P)
=\mathbf1\{\tau(P)>m_{0,\mathrm{len}}-a_{\mathrm{len}}\},
\qquad
D_{+,\mathrm{len}}(P)
=\mathbf1\{\tau(P)<m_{0,\mathrm{len}}+a_{\mathrm{len}}\}.
\]
On a left-target law, coverage implies
\(L_I(o)\le\tau(P)\le m_{0,\mathrm{len}}-a_{\mathrm{len}}\);
on a right-target law, it implies
\(m_{0,\mathrm{len}}+a_{\mathrm{len}}\le\tau(P)\le U_I(o)\).
Thus honesty from \cref{obj:def:interval-class} gives the respective
sample-event probabilities at least \(1-\alpha\).
On a failure law, \(1-\alpha-D_{\sigma,\mathrm{len}}(P)=-\alpha\le0\).
Consequently, for every \(P\in\mathcal M(d,q)\),
\[
\begin{aligned}
P_O^{\otimes n}(E_{\mathrm L,\alpha})
&\ge1-\alpha-D_{-,\mathrm{len}}(P),\\
P_O^{\otimes n}(E_{\mathrm R,\alpha})
&\ge1-\alpha-D_{+,\mathrm{len}}(P).
\end{aligned}
\]
Integrating and using the target-failure bounds gives
\[
\begin{aligned}
Q_{-,\mathrm{len}}^{(n)}(E_{\mathrm L,\alpha})
&\ge1-\alpha-\int D_{-,\mathrm{len}}\,d\Pi_{-,\mathrm{len}}
 \ge1-\alpha-\beta_\alpha,\\
Q_{+,\mathrm{len}}^{(n)}(E_{\mathrm R,\alpha})
&\ge1-\alpha-\int D_{+,\mathrm{len}}\,d\Pi_{+,\mathrm{len}}
 \ge1-\alpha-\beta_\alpha.
\end{aligned}
\]
The total-variation bound transfers the first estimate:
\[
Q_{+,\mathrm{len}}^{(n)}(E_{\mathrm L,\alpha})
\ge Q_{-,\mathrm{len}}^{(n)}(E_{\mathrm L,\alpha})-\beta_\alpha
\ge1-\alpha-2\beta_\alpha.
\]
Using the union-intersection identity and the fact that a union has
probability at most one, we obtain
\[
\begin{aligned}
Q_{+,\mathrm{len}}^{(n)}
 (E_{\mathrm L,\alpha}\cap E_{\mathrm R,\alpha})
&=
Q_{+,\mathrm{len}}^{(n)}(E_{\mathrm L,\alpha})
+Q_{+,\mathrm{len}}^{(n)}(E_{\mathrm R,\alpha})
-Q_{+,\mathrm{len}}^{(n)}
 (E_{\mathrm L,\alpha}\cup E_{\mathrm R,\alpha})\\
&\ge1-2\alpha-3\beta_\alpha.
\end{aligned}
\]
On the intersection, the endpoint inequalities imply
\[
U_I(o)-L_I(o)
\ge (m_{0,\mathrm{len}}+a_{\mathrm{len}})
   -(m_{0,\mathrm{len}}-a_{\mathrm{len}})
=2a_{\mathrm{len}}.
\]
Since interval length is nonnegative everywhere,
\[
2a_{\mathrm{len}}\,
\mathbf1_{E_{\mathrm L,\alpha}\cap E_{\mathrm R,\alpha}}(o)
\le U_I(o)-L_I(o).
\]
Integration under the plus mixture, followed by its defining mixture
identity, therefore gives
\[
\begin{aligned}
2a_{\mathrm{len}}\,
Q_{+,\mathrm{len}}^{(n)}
 (E_{\mathrm L,\alpha}\cap E_{\mathrm R,\alpha})
&\le
\int_{O^n}(U_I(o)-L_I(o))\,Q_{+,\mathrm{len}}^{(n)}(do)\\
&=
\int_{\mathcal M(d,q)}
\mathbb E_P[U_I(O^n)-L_I(O^n)]\,\Pi_{+,\mathrm{len}}(dP)\\
&\le
\sup_{P\in\mathcal M(d,q)}
\mathbb E_P[U_I(O^n)-L_I(O^n)].
\end{aligned}
\]
Multiplication of the intersection bound by \(2a_{\mathrm{len}}\ge0\)
and substitution of \(\beta_\alpha\) yield
\[
\begin{aligned}
\sup_{P\in\mathcal M(d,q)}
\mathbb E_P[U_I(O^n)-L_I(O^n)]
&\ge2a_{\mathrm{len}}(1-2\alpha-3\beta_\alpha)\\
&=2a_{\mathrm{len}}\frac58(1-2\alpha)\\
&=\frac54c_{\beta_\alpha}(1-2\alpha)g_{n,d,q}.
\end{aligned}
\]
Lengths lie in \([0,2]\), justifying these integrations, and the constant
interval \([-1,1]\) is honest because \(\tau(P)\in[-1,1]\).
Taking the infimum in \cref{obj:def:length-risk} proves
\[
\mathfrak L^*_{n,d,q,\alpha}
\ge c_{\mathrm{sep,len},\alpha}g_{n,d,q}
=\frac54c_{\beta_\alpha}(1-2\alpha)g_{n,d,q}.
\]
% lean: CausalSmith.Stat.MarNearcompleteFrontier.priorIntervalRisk_separation

\item We now combine the point-risk separation bound with the parametric
floor. By \cref{obj:def:rate},
\[
r_{n,d,q}=\frac1n+g_{n,d,q}^{\,2}.
\]
If \(g_{n,d,q}^{\,2}\ge n^{-1}\), then
\[
c\,r_{n,d,q}
\le 2c\,g_{n,d,q}^{\,2}
\le c_{\mathrm{sep,pt}}g_{n,d,q}^{\,2}
\le\mathfrak R^*_{n,d,q},
\]
where \(2c=\min\{c_{\mathrm{floor}},c_{\mathrm{sep,pt}}\}\le c_{\mathrm{sep,pt}}\).
If \(g_{n,d,q}^{\,2}<n^{-1}\), then
\[
c\,r_{n,d,q}
\le \frac{2c}{n}
\le \frac{c_{\mathrm{floor}}}{n}
\le\mathfrak R^*_{n,d,q},
\]
using \(2c\le c_{\mathrm{floor}}\) and the floor from
\cref{obj:lem:parametric-floor-infrastructure}.
Thus the combined point-risk bound holds throughout the parameter range.
% lean: CausalSmith.Stat.MarNearcompleteFrontier.pointRate_from_floor_and_separation

\item For the length bound, \(n\ge1\) ensures
\[
\sqrt n>0,
\qquad
(\sqrt n)^2=n,
\qquad
\frac1n=\left(\frac1{\sqrt n}\right)^2.
\]
Also \(r_{n,d,q}\ge0\), so
\[
\sqrt{r_{n,d,q}}\ge0,
\qquad
\bigl(\sqrt{r_{n,d,q}}\bigr)^2=r_{n,d,q}.
\]
If \(g_{n,d,q}^{\,2}\ge n^{-1}\), then
\[
\bigl(\sqrt{r_{n,d,q}}\bigr)^2
=\frac1n+g_{n,d,q}^{\,2}
\le2g_{n,d,q}^{\,2}
\le(2g_{n,d,q})^2.
\]
Both square roots being compared are nonnegative, hence
\[
\sqrt{r_{n,d,q}}\le2g_{n,d,q}.
\]
It follows that
\[
c_\alpha\sqrt{r_{n,d,q}}
\le2c_\alpha g_{n,d,q}
\le c_{\mathrm{sep,len},\alpha}g_{n,d,q}
\le\mathfrak L^*_{n,d,q,\alpha}.
\]
If \(g_{n,d,q}^{\,2}<n^{-1}\), then
\[
\bigl(\sqrt{r_{n,d,q}}\bigr)^2
=\frac1n+g_{n,d,q}^{\,2}
\le\frac2n
\le\left(\frac2{\sqrt n}\right)^2,
\]
and nonnegativity gives
\[
\sqrt{r_{n,d,q}}\le\frac2{\sqrt n}.
\]
Therefore
\[
c_\alpha\sqrt{r_{n,d,q}}
\le\frac{2c_\alpha}{\sqrt n}
=\frac{\min\{c_{0,\alpha},c_{\mathrm{sep,len},\alpha}\}}{\sqrt n}
\le\frac{c_{0,\alpha}}{\sqrt n}
\le\mathfrak L^*_{n,d,q,\alpha},
\]
using \cref{obj:lem:parametric-floor-infrastructure}.
These two cases cover every admissible parameter choice, including
\(q=1\), for which \(g_{n,d,q}=0\) and the second case applies.
Together with the separation bounds and the parametric floors established
above, this proves all the asserted guarantees.
% lean: CausalSmith.Stat.MarNearcompleteFrontier.intervalRate_from_floor_and_separation
\end{enumerate}
\end{proof}
